\begin{table}[!htbp]
\centering
\begin{ruledtabular}
\begin{tabular}{ccccc}
$\Phi$ & $C$ & full coefficient & profile/scalar coefficient & relative deviation\\
\colrule
$1.000$ & $0.000000$ & $0.000000$ & $0.000000$ & ---\\
$0.950$ & $1.453256$ & $1.906955$ & $1.931512$ & $1.288\%$\\
$0.800$ & $3.316601$ & $4.368046$ & $4.408069$ & $0.916\%$\\
$0.600$ & $4.992296$ & $6.593434$ & $6.635222$ & $0.634\%$\\
$0.500$ & $5.645751$ & $7.465353$ & $7.503724$ & $0.514\%$\\
$0.400$ & $6.192296$ & $8.197671$ & $8.230133$ & $0.396\%$\\
$0.200$ & $6.916601$ & $9.180756$ & $9.192801$ & $0.131\%$\\
$0.125$ & $7.000000$ & $9.303646$ & $9.303646$ & $0.000\%$\\
\end{tabular}
\end{ruledtabular}
\caption{Boundary scan for the same $d=8$ mesh as Table~\ref{tab:reck}. Entries are coefficients of $\sigma^2$. The $\Phi=0.50$ row is also the quadratic prediction in Fig.~\ref{fig:quadraticscaling}. The undefined relative deviation at zero loss is left blank; no numerical quantity is missing.}
\label{tab:d8scan}
\end{table}
